\documentclass[11pt]{article}

%\renewcommand\familydefault{\sfdefault}
 
%input preamble and macros
\input{preamble}
\input{macros}

 
\usepackage[margin=1in]{geometry} 
\usepackage{amsmath,amsthm,amssymb}
\usepackage{hyperref}
\usepackage[linewidth=1pt]{mdframed}
%\newcommand{\N}{\mathbb{N}}
%\newcommand{\Z}{\mathbb{Z}}

% Named environments (no counters) 
\newenvironment{theorem}[2][Theorem]{\begin{trivlist}
\item[\hskip \labelsep {\bfseries #1}\hskip \labelsep {\bfseries #2.}]}{\end{trivlist}}
\newenvironment{lemma}[2][Lemma]{\begin{trivlist}
\item[\hskip \labelsep {\bfseries #1}\hskip \labelsep {\bfseries #2.}]}{\end{trivlist}}
%\newenvironment{exercise}[2][Exercise]{\begin{trivlist}
%\item[\hskip \labelsep {\bfseries #1}\hskip \labelsep {\bfseries #2.}]}{\end{trivlist}}
\newenvironment{problem}[2][Problem]{\begin{trivlist}
\item[\hskip \labelsep {\bfseries #1}\hskip \labelsep {\bfseries #2.}]}{\end{trivlist}}
\newenvironment{question}[2][Question]{\begin{trivlist}
\item[\hskip \labelsep {\bfseries #1}\hskip \labelsep {\bfseries #2.}]}{\end{trivlist}}
\newenvironment{corollary}[2][Corollary]{\begin{trivlist}
\item[\hskip \labelsep {\bfseries #1}\hskip \labelsep {\bfseries #2.}]}{\end{trivlist}}
 
% Environments with counters
\newtheoremstyle{myplain} {5mm}% ⟨Space above⟩
{3mm}% ⟨Space below⟩
{}% ⟨Body font⟩
{}% ⟨Indent amount⟩
{\bfseries}% ⟨Theorem head font⟩
{.}% ⟨Punctuation after theorem head⟩
{.5em}% ⟨Space after theorem head⟩2
{}% ⟨Theorem head spec (can be left empty, meaning ‘normal’)⟩

\date{}

\theoremstyle{myplain}
\newtheorem{exercise}{Exercise}

\theoremstyle{definition} % no italics
\newtheorem{subexercise}{}[exercise]

\newcommand{\subex}[1]{\begin{subexercise}#1\end{subexercise}}

\begin{document}
 
% --------------------------------------------------------------
%                         Start here
% --------------------------------------------------------------
 
\title{Modelling and analysis of a cyber-physical system \\
now with \underline{monads} !! \\[10pt] \small{Practical assignment 2}}
\author{Renato Neves}
 
\maketitle

\noindent
In the middle of the night four adventurers encounter a rope-bridge that spans
a deep ravine. For safety reasons, no more than two people cross the bridge at
the same time and a flashlight needs to be carried in every crossing.  There is
only one flashlight. The four adventurers are not equally skilled: crossing the
bridge takes them one, two, five, and ten minutes respectively. A pair of
adventurers crosses the bridge in an amount of time equal to that of the
slowest of the two.
\begin{center}
\includegraphics[scale=0.8]{images/ropebridge-diag}
\end{center}
One adventurer claims that they cannot be all on the other side in less than 19
minutes. Another one disagrees and claims that it can be done in 17 minutes.

\subsubsection*{First task}

Your first task is to verify these claims using
\textsc{Haskell}\footnote{An animated description of the problem is
  available
  \href{https://www.youtube.com/watch?v=7yDmGnA8Hw0}{here}.}. Specifically
you will need to,
\begin{enumerate}
\item \underline{model} the system above using what you learned about
  \underline{monads}, in particular the duration and non-deterministic ones;
\item show that it is indeed \underline{possible} for all adventurers
  to be on the other side in 17 minutes;
\item show that it is \underline{impossible} for all adventurers
  to be on the other side in less than 17 minutes.
\end{enumerate}
You should fulfill this task by \underline{completing the code} in attachment
(\texttt{Adventurers.hs}), i.e. by adding a definition to the functions that
lack a definition, \underline{following the comments} present in the code. You
are free to add new auxiliary functions.

\noindent
\textbf{Some hints} to help you get started: Recall the duration monad from the
slides and the \textsc{Haskell} code that was previously provided. Analyse in
detail the code concerning the Knight's quest and in particular the monad
\texttt{LogList}.

\subsubsection*{Second task}

Implement a function that receives a list of valid plays and returns a
\emph{distribution} of the resulting states together with the corresponding
elapsed time. The function should in particular encompass \emph{probabilistic
crossing times}: \ie\ if the adventurer's crossing time is $x$ then its
probabilistic crossing time will now be given by the uniform distribution,
\[
        \frac{1}{3} \cdot \left ( x - \frac{x}{2} \right ) 
                + \frac{1}{3} \cdot  \left ( x \right )
                + \frac{1}{3} \cdot \left ( x + \frac{x}{2} \right )
\]
In this probabilistic setting, what can you conclude about the adventurer's
claims mentioned above?  Again you should fulfill this task by
\underline{completing the code} in attachment (\texttt{Adventurers.hs}), \ie\
by adding a definition to the functions that lack a definition,
\underline{following the comments} present in the code. Feel free to add new
auxiliary functions.

\subsubsection*{Third task}

Add new functionalities to the code which you think are essential or
interesting. For example present the sequence of states respective to the
movements of the adventurers from the initial state to the final goal, or make
the game more interactive via the IO monad. Imagination is the limit -- but of
course remind yourself that the focus is the use of monads in programming :-).

\bigskip

\begin{mdframed}  
  \myparagraph{What to submit:} A single report in PDF for all tasks and the
  respective source files. Send by email (\underline{nevrenato@gmail.com}) a
  unique zip file ``\bash{cf2526-N.zip}'', where \bash{N} is your student
  number. The subject of the email should be ``\bash{cf2526 N}''.

\myparagraph{Deadline:} 10th June 2026 @ 23h59
\end{mdframed}

\end{document}
